\documentclass[11pt,fontset=none]{ctexart}

\usepackage[a4paper,margin=1in]{geometry}
\usepackage{amsmath,amssymb}
\usepackage{booktabs}
\usepackage{array}
\usepackage{graphicx}
\usepackage{float}
\usepackage{hyperref}
\usepackage{tikz}
\usepackage{xcolor}
\usetikzlibrary{arrows.meta,positioning,calc,fit}

\setCJKmainfont{Songti SC}
\setCJKsansfont{Heiti SC}

\hypersetup{
  colorlinks=true,
  linkcolor=blue!45!black,
  urlcolor=blue!45!black,
  citecolor=blue!45!black
}

\title{ARIA-GOLF v1--v5：结构差异与模型公式}
\author{内部架构说明}
\date{\today}

\newcommand{\fs}{f_{\mathrm{s}}}
\newcommand{\src}{s}
\newcommand{\harm}{h}
\newcommand{\noise}{\eta}
\newcommand{\vt}{H_{\mathrm{VT}}}
\newcommand{\room}{H_{\mathrm{room}}}
\newcommand{\nfilt}{H_{\mathrm{N}}}
\newcommand{\ap}{A}
\newcommand{\lf}{g_{\mathrm{LF}}}
\newcommand{\biquad}{B}
\newcommand{\pole}{P}
\newcolumntype{L}[1]{>{\raggedright\arraybackslash}p{#1}}

\begin{document}
\maketitle

\begin{abstract}
本文把 ARIA-GOLF v1--v5 的配置差异整理成一个便于理解和展示的模型比较说明。
五个版本共享相同的编码器系列、观测到的 $F_0$ 条件、LF 声门流振荡器、白噪声发生器、
room filter 以及 multi-scale spectrogram loss。真正的结构差异集中在两个位置：
一是 v1--v3 逐步调整的解析声道级联滤波器，二是 v3--v5 改动的谐波/噪声激励拓扑。
在所有版本中，谐波支路和噪声支路先在 source level 合成为一个激励信号 $\src$，
然后这个合成后的 source 再整体经过一次 VocalTractCascade / LPC-like 声道滤波器。
下面的公式把这些变化写成一个可解释的 ablation sequence。
\end{abstract}

\section*{如何阅读这五个版本}

最简单的读法是：v1--v5 不是五个互不相关的系统，而是沿着两条轴线展开的结构消融。
第一条轴线是 v1--v3：它们主要在调整解析 vocal tract，即显式 formant control、
bandwidth constraint、spectral tilt 和 learned residual all-pole capacity 之间的平衡。
第二条轴线是 v3--v5：v4 和 v5 保留 v3 的 vocal tract，只改变 excitation mixture，
因为 v3 剩下的主要问题来自 additive harmonic-plus-noise source 中的高频周期性条纹。

\begin{figure}[H]
\centering
\resizebox{\linewidth}{!}{%
\begin{tikzpicture}[
  font=\sffamily\small,
  >=Latex,
  model/.style={draw=black, rounded corners=2pt, align=left, text width=2.9cm,
    minimum height=1.55cm, inner sep=5pt, fill=white},
  vt/.style={model, fill=blue!5},
  src/.style={model, fill=orange!10},
  band/.style={draw=black!40, dashed, rounded corners=3pt, inner sep=7pt},
  note/.style={align=center, font=\sffamily\footnotesize},
  arrow/.style={->, line width=0.8pt}
]
  \node[vt] (v1) at (0,0) {\textbf{v1}\\tilt on\\宽 bandwidth\\无 F supervision};
  \node[vt] (v2) at (4.15,0) {\textbf{v2}\\tilt off\\紧 bandwidth\\additive source};
  \node[vt] (v3) at (8.3,0) {\textbf{v3}\\tilt restored\\紧 bandwidth\\additive source};
  \node[src] (v4) at (13.05,1.18) {\textbf{v4}\\v3 vocal tract\\独立 AP head $A$};
  \node[src] (v5) at (13.05,-1.18) {\textbf{v5}\\v3 vocal tract\\bounded crossfade $q$};

  \draw[arrow] (v1) -- (v2);
  \draw[arrow] (v2) -- (v3);
  \coordinate (split) at ($(v3.east)+(0.9,0)$);
  \draw[arrow] (v3.east) -- (split) |- (v4.west);
  \draw[arrow] (v3.east) -- (split) |- (v5.west);
  \node[band, fit=(v1)(v2)(v3)] (vtband) {};
  \node[band, draw=orange!70!black, fit=(v4)(v5)] (srcband) {};
  \node[note, blue!55!black, above=0.08cm of vtband] {vocal-tract refinements};
  \node[note, orange!75!black, above=0.08cm of srcband] {excitation fixes};
\end{tikzpicture}}
\vspace{2pt}

{\footnotesize v1$\rightarrow$v2：收紧 F/B 范围 \qquad
v2$\rightarrow$v3：恢复 spectral tilt}
\caption{ARIA-GOLF v1--v5 的概念谱系。v1--v3 应读作声道滤波器的逐步改进；
v4 和 v5 是针对同一个高频 excitation artifact 的两种修复方案。v1 同时也是
没有显式 formant supervision 的 baseline。}
\label{fig:lineage-zh}
\end{figure}

\begin{figure}[H]
\centering
\resizebox{\linewidth}{!}{%
\begin{tikzpicture}[
  font=\sffamily\small,
  >=Latex,
  block/.style={draw=black, rounded corners=2pt, align=center,
    minimum width=1.45cm, minimum height=0.72cm, inner sep=3pt, fill=white},
  sum/.style={draw=black, circle, align=center, minimum size=0.65cm, inner sep=0pt, fill=white},
  panel/.style={draw=black!45, rounded corners=3pt, inner sep=7pt},
  flow/.style={->, line width=0.8pt},
  lbl/.style={font=\sffamily\footnotesize, align=center}
]
  \begin{scope}[xshift=0cm]
    \node[lbl] at (2.2,2.2) {\textbf{v1--v3: additive}\\只能加噪声，不能减谐波};
    \node (h1) at (0,1.25) {$h$};
    \node (n1) at (0,0) {$\eta$};
    \node[block] (hn1) at (1.75,0) {$H_N^{\exp}$};
    \node[sum] (sum1) at (3.35,0.65) {$+$};
    \node (s1) at (4.45,0.65) {$s$};
    \draw[flow] (h1) -- (sum1);
    \draw[flow] (n1) -- (hn1);
    \draw[flow] (hn1) -- (sum1);
    \draw[flow] (sum1) -- (s1);
    \node[lbl] at (2.2,-0.9) {$s=h+H_N\eta$};
    \node[panel, fit=(h1)(n1)(hn1)(sum1)(s1)] {};
  \end{scope}

  \begin{scope}[xshift=5.5cm]
    \node[lbl] at (2.4,2.2) {\textbf{v4: decoupled AP}\\噪声颜色与混合权重分离};
    \node (h2) at (0,1.25) {$h$};
    \node (n2) at (0,0) {$\eta$};
    \node[block] (onea) at (1.55,1.25) {$1-A$};
    \node[block] (hn2) at (1.55,0) {$H_N^{\exp}$};
    \node[block] (a2) at (3.0,0) {$A$};
    \node[sum] (sum2) at (4.25,0.65) {$+$};
    \node (s2) at (5.25,0.65) {$s$};
    \draw[flow] (h2) -- (onea);
    \draw[flow] (onea) -- (sum2);
    \draw[flow] (n2) -- (hn2);
    \draw[flow] (hn2) -- (a2);
    \draw[flow] (a2) -- (sum2);
    \draw[flow] (sum2) -- (s2);
    \draw[dashed, line width=0.6pt] (a2.north) -- (onea.south);
    \node[lbl] at (2.65,-0.9) {$S=(1-A)Y+A\,H_NN$};
    \node[panel, fit=(h2)(n2)(onea)(hn2)(a2)(sum2)(s2)] {};
  \end{scope}

  \begin{scope}[xshift=12.2cm]
    \node[lbl] at (2.25,2.2) {\textbf{v5: bounded crossfade}\\一个滤波器承担两个职责};
    \node (h3) at (0,1.25) {$h$};
    \node (n3) at (0,0) {$\eta$};
    \node[block] (oneq) at (1.65,1.25) {$1-q$};
    \node[block] (q3) at (1.65,0) {$q=\sigma(r)$};
    \node[sum] (sum3) at (3.35,0.65) {$+$};
    \node (s3) at (4.45,0.65) {$s$};
    \draw[flow] (h3) -- (oneq);
    \draw[flow] (oneq) -- (sum3);
    \draw[flow] (n3) -- (q3);
    \draw[flow] (q3) -- (sum3);
    \draw[flow] (sum3) -- (s3);
    \draw[dashed, line width=0.6pt] (q3.north) -- (oneq.south);
    \node[lbl] at (2.2,-0.9) {$S=(1-q)Y+qN$};
    \node[panel, fit=(h3)(n3)(oneq)(q3)(sum3)(s3)] {};
  \end{scope}
\end{tikzpicture}}
\caption{三种 excitation topology。Additive synthesis 可以增加噪声地板，但不能抑制已经存在的周期性谐波。
v4 和 v5 引入 per-band harmonic/noise mixture，因此可以把一部分高频周期性能量转成 shaped noise。}
\label{fig:excitation-zh}
\end{figure}

\section{共享的 Source-Filter 骨架}

令 $x$ 表示输入波形，$F_{0,k}$ 表示观测到的 frame-rate 基频轨迹。编码器预测 frame-rate control streams：
\begin{equation}
  \mathbf{c}_k = E_\theta(x, F_{0,k}),
\end{equation}
其中最后一层 projection 的宽度由 decoder 声明的 control splits 决定。上采样之后，LF source 的相位由下式得到：
\begin{equation}
  \phi_n = \phi_{n-1} + 2\pi \frac{F_{0,n}}{\fs},
\end{equation}
deterministic harmonic source 为：
\begin{equation}
  \harm_n = v_n \, \lf(\phi_n; R_{d,n}),
\end{equation}
其中 $v_n=\mathbb{1}[F_{0,n}>0]$ 是 voicing gate，$R_d$ 控制 LF glottal-flow shape。
每个版本都有一个白噪声激励 $\noise_n\sim\mathcal{N}(0,1)$。模型首先把 harmonic branch
和 noise branch 合成为一个 source/excitation signal $\src^{(m)}$，然后最终波形总是：
\begin{equation}
  \hat{x} = \room * \bigl(\vt^{(m)} * \src^{(m)}\bigr),
\end{equation}
其中 $m\in\{1,\ldots,5\}$ 表示模型版本。因此，各版本只在 version-specific source
$\src^{(m)}$ 和 vocal-tract filter $\vt^{(m)}$ 上不同。实现上，这是对合成后的 source
做一次 vocal-tract filtering，也就是 $\mathrm{LPC}(\harm+\mathrm{noise})$，
而不是先分别滤波再相加。由于 end filter 是线性的，两种写法在数学上等价；
但从架构和代码顺序看，source-level 合成再过 vocal tract 是更清楚的读法。

\section{Vocal-Tract Cascade}

解析 vocal tract 包含一个 gain term、一个可选 spectral-tilt term、两个 deterministic
all-pole formant sections，以及若干 learned biquad sections：
\begin{equation}
  {\vt}^{(m)}_n(z)
  =
  G_n
  \, T_{\alpha,n}^{\delta_m}(z)
  \prod_{j=1}^{2} \pole(z; F_{j,n}, B_{j,n})
  \prod_{\ell=1}^{L_m} \biquad_{\ell,n}(z).
  \label{eq:vt-general-zh}
\end{equation}
这里 $\delta_m=1$ 表示启用 spectral tilt，$\delta_m=0$ 表示不启用。Learned residual
sections $\biquad_{\ell,n}$ 用来吸收显式 F1/F2 poles 没有覆盖到的 vocal-tract detail。
下表报告的是 16 kHz F024 配置，也就是 v4/v5 所在的配置。

\begin{table}[H]
\centering
\small
\begin{tabular}{@{}lcccccc@{}}
\toprule
版本 & Tilt $\delta_m$ & Learned $L_m$ & F1 range & F2 range & B1/B2 range & LPC order \\
\midrule
v1 & 1 & 6 & 150--1300 & 600--3200 & 30--330 / 30--430 & 18 \\
v2 & 0 & 7 & 200--1000 & 700--3000 & 50--200 / 50--250 & 18 \\
v3 & 1 & 6 & 150--1300 & 600--3200 & 50--200 / 50--250 & 18 \\
v4 & 1 & 6 & 150--1300 & 600--3200 & 50--200 / 50--250 & 18 \\
v5 & 1 & 6 & 150--1300 & 600--3200 & 50--200 / 50--250 & 18 \\
\bottomrule
\end{tabular}
\caption{16 kHz F024 配置下 v1--v5 的 vocal-tract 差异。v4 和 v5 继承 v3 的
vocal-tract design，它们的变化发生在 excitation topology。}
\end{table}

这里的 nominal LPC order 按 biquad-like sections 的数量计算：
\begin{equation}
  O_m = 2\bigl(\delta_m + 2 + L_m\bigr).
\end{equation}
因此 16 kHz 下 v1/v3/v4/v5 的 vocal tract 是 $2(1+2+6)=18$，而 v2 是
$2(0+2+7)=18$。Tilt section $T_{\alpha}(z)$ 是一个 degenerate biquad，
只有一个 real pole，所以旧注释里可能会把 tilt-enabled 16 kHz model 写成 effective
pole count 17。本文统一使用 nominal section-count convention。

\begin{table}[H]
\centering
\small
\begin{tabular}{@{}lcccc@{}}
\toprule
采样率 / speaker & Nominal LPC order & v1 $L_m$ & v2 $L_m$ & v3 $L_m$ \\
\midrule
16 kHz / F024 & 18 & 6 & 7 & 6 \\
22 kHz / LJSpeech & 22 & 8 & 9 & 8 \\
24 kHz / CSMSC & 22 & 8 & 9 & 8 \\
\bottomrule
\end{tabular}
\caption{Learned residual cascade 的采样率依赖性。解析 F1/F2 controls 和 spectral tilt
在概念上保持不变；更高采样率使用更多 learned biquads 来覆盖额外的高频细节。}
\end{table}

\section{各版本公式}

\subsection{v1：第一个 Analytic-Formant Baseline}

v1 引入了带 spectral tilt 的可解释 F1/F2 vocal-tract cascade，bandwidth 约束较宽，
并包含 6 个 learned residual biquads：
\begin{align}
  {\vt}^{(1)}_n(z)
  &=
  G_n T_{\alpha,n}(z)
  \prod_{j=1}^{2}\pole(z;F_{j,n},B_{j,n})
  \prod_{\ell=1}^{6}\biquad_{\ell,n}(z),\\
  \src^{(1)}_n
  &=
  \harm_n + ({\nfilt}^{\exp}_n * \noise)_n.
\end{align}
Additive excitation 可以抬高噪声地板，但不能削弱 $\harm_n$ 中已经存在的 coherent
high-frequency harmonics。需要特别注意的是，F024 训练 recipe 中 v1 的 F1/F2 control
没有显式 formant supervision，它们只通过 reconstruction objective 间接形成，因此 v1
是 formant control 最松的 baseline。

\subsection{v2：无 Spectral Tilt 的更强 Formant Control}

v2 去掉 tilt branch，收紧 F1/F2 和 B1/B2 的范围，并增加一个 learned biquad：
\begin{align}
  {\vt}^{(2)}_n(z)
  &=
  G_n
  \prod_{j=1}^{2}\pole(z;F_{j,n},B_{j,n})
  \prod_{\ell=1}^{7}\biquad_{\ell,n}(z),\\
  \src^{(2)}_n
  &=
  \harm_n + ({\nfilt}^{\exp}_n * \noise)_n.
\end{align}
v2 的目的在于增强 formant manipulability；但去掉 tilt 后，broad spectral slope
就必须更多地由 residual sections 去解释。

\subsection{v3：当前 Vocal-Tract Base}

v3 恢复 spectral tilt，同时保留 v2 中更紧的 bandwidth constraints，并回到 6 个
learned biquads：
\begin{align}
  {\vt}^{(3)}_n(z)
  &=
  G_n T_{\alpha,n}(z)
  \prod_{j=1}^{2}\pole(z;F_{j,n},B_{j,n})
  \prod_{\ell=1}^{6}\biquad_{\ell,n}(z),\\
  \src^{(3)}_n
  &=
  \harm_n + ({\nfilt}^{\exp}_n * \noise)_n.
\end{align}
这是最强的 analytic vocal-tract baseline：它同时包含 F1/F2 control、bounded
bandwidths、spectral tilt 和 residual all-pole capacity。它的主要弱点仍然是
additive-only excitation，因此在强 voiced region 里可能留下 high-frequency harmonic stripes。

\subsection{v4：Decoupled Band-Aperiodicity Source}

v4 保留 v3 的 vocal tract，但加入一个显式 aperiodicity filter
$\ap_n(\omega)\in[0,1]$，并带有 rising high-frequency prior：
\begin{align}
  {\vt}^{(4)}_n(z) &= {\vt}^{(3)}_n(z),\\
  \ap_n(\omega)
  &=
  \sigma\bigl(a_n(\omega)+b_{\mathrm{HF}}(\omega)\bigr),\\
  S^{(4)}_n(\omega)
  &=
  \bigl(1-\ap_n(\omega)\bigr)Y_n(\omega)
  +
  \ap_n(\omega){\nfilt}^{\exp}_n(\omega)N_n(\omega).
\end{align}
这里 $Y_n(\omega)$ 和 $N_n(\omega)$ 分别是 harmonic source 和 white noise 的 short-time
spectra。Noise-colour filter $\nfilt$ 保持 unbounded，而 $\ap$ 单独控制 harmonic/noise
mixture。因此 v4 是修复 high-band harmonic-stripe artifact 最清晰、最原则化的版本。

\subsection{v5：Minimal Bounded Crossfade}

v5 也保留 v3 的 vocal tract，但把 noise filter 本身用作 bounded crossfade：
\begin{align}
  {\vt}^{(5)}_n(z) &= {\vt}^{(3)}_n(z),\\
  q_n(\omega) &= \sigma(r_n(\omega)),\\
  S^{(5)}_n(\omega)
  &=
  \bigl(1-q_n(\omega)\bigr)Y_n(\omega)
  +
  q_n(\omega)N_n(\omega).
\end{align}
等价地，用 operator form 可以写成：
\begin{equation}
  \src^{(5)}
  =
  (I-Q)\harm + Q\noise
  =
  \harm + Q\noise - Q\harm.
\end{equation}
这是最小结构改动：不增加额外 aperiodicity head，但同一个 bounded filter 必须同时承担
noise colour 和 aperiodicity weight 两个职责。

\section{作为 Ablation Sequence 的解释}

\begin{table}[H]
\centering
\small
\begin{tabular}{@{}lccccc@{}}
\toprule
版本 & Formant loss & Smoothness & Residual reg. & Batch $\times$ steps & Data passes \\
\midrule
v1 & 0 & -- & -- & $64\times40$k & $1\times$ \\
v2 & 1.5 & -- & 0.02 & $64\times80$k & $2\times$ \\
v3 & 2.0 & 0.05 & 0.02 & $128\times40$k & $2\times$ \\
v4 & 2.0 & 0.05 & 0.02 & $128\times40$k & $2\times$ \\
v5 & 2.0 & 0.05 & 0.02 & $128\times40$k & $2\times$ \\
\bottomrule
\end{tabular}
\caption{F024 比较中的训练和 supervision 差异。关键修正是：v1 没有显式 formant supervision；
不应再使用 ``light supervision'' 的说法。}
\end{table}

\begin{table}[H]
\centering
\small
\begin{tabular}{@{}L{0.09\linewidth}L{0.23\linewidth}L{0.34\linewidth}L{0.24\linewidth}@{}}
\toprule
版本 & 意图 & 主要变化 & 代价 / 风险 \\
\midrule
v1 & 第一个 analytic-formant ARIA & Tilt on，宽 bandwidth，无 formant supervision & 重建尚可，但 formant control 松散、约束较弱 \\
v2 & 更强 formant control & Tilt off，收紧 bandwidth，七个 learned sections & Control 改善，但 spectral slope 不够显式 \\
v3 & 当前 base & Tilt restored，紧 bandwidth，强 supervision & 操控较好，但 additive source artifact 仍在 \\
v4 & 原则化 artifact fix & 加入 decoupled band-aperiodicity $\ap$ & 表达力最强，但多一个 control head \\
v5 & 最小 artifact fix & Bounded noise filter 变成 crossfade $q$ & 不增加 head，但 noise level 与 aperiodicity 耦合 \\
\bottomrule
\end{tabular}
\caption{v1--v5 的推荐叙事方式。}
\end{table}

\section{评估含义}

v1--v3 的比较应主要关注 formant-control error、copy-synthesis quality，以及 spectral tilt
是否能提升自然度而不削弱操控。v3--v5 的比较则应聚焦 high-frequency voiced regions，
因为核心假设是：additive excitation 无法移除 coherent high-band harmonics，而 v4 和 v5
可以把一部分 harmonic source 转换为 shaped noise。最有信息量的可视化包括：
\begin{enumerate}
  \item 一张 lineage diagram，把 v1--v3 显示为 vocal-tract refinements，
        把 v4--v5 显示为 excitation fixes；
  \item 一张 side-by-side excitation-flow figure，对比 additive、decoupled aperiodicity
        和 bounded crossfade；
  \item 一个 spectrogram strip，对比 v3、v4、v5 在 4--8 kHz 区间的表现；
  \item 在 voiced frames 上展示 v4 的 $\ap_n(\omega)$ 和 v5 的 $q_n(\omega)$。
\end{enumerate}

\end{document}
